\documentclass[troisieme]{fiche}
\metadonnees{18}{Le prince et l'hirondelle}{Bloc D — Rapporter et supposer}

\begin{document}
\entete

\objectifs{%
\textbf{Langue} : l'hypothèse réelle, \emph{if} + présent ; \emph{when}, \emph{unless},
\emph{as soon as}.\par
\textbf{Texte} : Oscar Wilde, \og The Happy Prince \fg{} (1888).\par
\textbf{Durée} : 1 h — exercices 20 min, texte et questions 25 min, version 15 min.}

%% ===================================================================
\exercice[12 min]{\emph{If} + présent}

\rappel{\textbf{Rappel.} Après \emph{if}, jamais de \emph{will} : \emph{if it rains, I will
stay}. Le présent anglais y porte le futur.}

\consigne{Mettez le verbe entre parenthèses au temps qui convient.}

\begin{anglais}
\begin{enumerate}[label=\arabic*.,itemsep=2.5mm,leftmargin=*]
  \item If it \emph{(rain)} \trou{rains} tonight, the Swallow will be wet again.
  \item The boy will die if nobody \emph{(help)} \trou{helps} him.
  \item If you \emph{(stay)} \trou{stay} one night, you \emph{(be)} \trou{will be} my
    messenger.
  \item If I \emph{(go)} \trou{go} to Egypt now, I will never see him again.
  \item When the sun \emph{(rise)} \trou{rises}, the Swallow will fly away.
  \item As soon as he \emph{(arrive)} \trou{arrives}, he will tell the Prince.
  \item Unless you \emph{(leave)} \trou{leave} tonight, the winter \emph{(catch)}
    \trou{will catch} you.
  \item I will not go if you \emph{(ask)} \trou{ask} me to stay.
  \item If the statue \emph{(not be)} \trou{is not} solid gold, the Mayor will be angry.
  \item I don't know if he \emph{(come)} \trou{will come} back.
\end{enumerate}
\end{anglais}

\note[La règle et sa raison]{Après \emph{if}, \emph{when}, \emph{as soon as}, \emph{until},
\emph{unless}, l'anglais emploie le présent pour parler de l'avenir : la condition est posée
comme un fait, c'est la conséquence qui reste à venir. Le français, lui, met les deux au futur
(\og quand il fera nuit \fg).}

\note[Les deux \emph{if}]{Phrase 10 : \emph{if} y signifie \og si \fg{} au sens de \og
est-ce que \fg, il introduit une question rapportée et non une condition. \emph{Will} y est
donc normal. C'est le seul cas.}

\note[\emph{Unless}]{= \emph{if\ldots not}. \emph{Unless you leave} = \emph{if you don't
leave}. Il ne se met donc jamais à la forme négative.}

%% ===================================================================
\exercice[8 min]{Traduire la condition}

\rappel{\textbf{Rappel.} \og Quand il fera \fg, \og dès qu'elle arrivera \fg{} : présent en
anglais.}

\consigne{Traduisez en anglais.}

\begin{anglais}
\begin{enumerate}[label=\arabic*.,itemsep=3mm,leftmargin=*]
  \item Si tu pars, je resterai seul.
    \lignes{1}
    \corr{If you leave, I will be alone.}
  \item Quand il fera nuit, il s'envolera.
    \lignes{1}
    \corr{When it is dark, he will fly away.}
  \item Je ne partirai pas tant qu'il ne sera pas guéri.
    \lignes{1}
    \corr{I will not leave until he is better.}
  \item À moins qu'il ne pleuve, nous sortirons.
    \lignes{1}
    \corr{Unless it rains, we will go out.}
  \item Dès qu'elle arrivera, préviens-moi.
    \lignes{1}
    \corr{As soon as she arrives, tell me.}
  \item Je ne sais pas s'il viendra.
    \lignes{1}
    \corr{I don't know if he will come.}
\end{enumerate}
\end{anglais}

\note[Deux pièges de traduction]{\og Tant que\ldots ne pas \fg{} se dit \emph{until} +
affirmative : \emph{until he is better}, et non \emph{until he is not better}. \og À moins
que \fg{} se dit \emph{unless}, sans négation non plus.}

%% ===================================================================
\newpage
\section{Texte}

\noindent\small\textit{Une statue dorée domine la ville. Une hirondelle en retard sur ses
compagnes, en route pour l'Égypte, s'y arrête pour la nuit.}\normalsize

\begin{texte}
``What is the use of a statue if it cannot keep the rain off?'' he said; ``I must look for a
good \gl[a chimney-pot]{chimney-pot}{une cheminée (le conduit)},'' and he determined to fly
away.

But before he had opened his wings, a third drop fell, and he looked up, and saw --- Ah! what
did he see?

The eyes of the Happy Prince were filled with tears, and tears were running down his golden
\gl[a cheek]{cheeks}{une joue}. His face was so beautiful in the moonlight that the little
Swallow was filled with pity.

``Who are you?'' he said.

``I am the Happy Prince.''

``Why are you \gl[to weep]{weeping}{pleurer} then?'' asked the Swallow; ``you have quite
\gl[to drench]{drenched}{tremper} me.''

``When I was alive and had a human heart,'' answered the statue, ``I did not know what tears
were, for I lived in the Palace of Sans-Souci, where \gl[sorrow]{sorrow}{le chagrin} is not
allowed to enter. In the daytime I played with my companions in the garden, and in the evening
I led the dance in the Great Hall. Round the garden ran a very \gl[lofty]{lofty}{très haut}
wall, but I never cared to ask what lay beyond it, everything about me was so beautiful. My
\gl[a courtier]{courtiers}{un courtisan} called me the Happy Prince. So I lived, and so I
died. And now that I am dead they have set me up here so high that I can see all the ugliness
and all the misery of my city, and though my heart is made of \gl[lead]{lead}{le plomb} yet I
cannot choose but weep.''

``What! is he not solid gold?'' said the Swallow to himself. He was too polite to make any
personal remarks out loud.

``Far away,'' continued the statue in a low musical voice, ``far away in a little street there
is a poor house. One of the windows is open, and through it I can see a woman
\gl[to seat]{seated}{assis} at a table. Her face is thin and \gl[worn]{worn}{usé}, and she has
\gl[coarse]{coarse}{rêche, grossier}, red hands, all \gl[to prick]{pricked}{piquer} by the
\gl[a needle]{needle}{une aiguille}, for she is a
\gl[a seamstress]{seamstress}{une couturière}. She is
\gl[to embroider]{embroidering}{broder} passion-flowers on a satin gown for the
\gl[loveliest]{loveliest}{la plus ravissante} of the Queen's maids-of-honour to wear at the
next Court-ball. In a bed in the corner of the room her little boy is lying ill. He has a
fever, and is asking for oranges. His mother has nothing to give him but river water, so he
is crying. Swallow, Swallow, little Swallow, will you not bring her the ruby out of my
\gl[a sword-hilt]{sword-hilt}{une poignée d'épée}? My feet are
\gl[to fasten]{fastened}{attacher} to this \gl[a pedestal]{pedestal}{un socle} and I cannot
move.''

``I am waited for in Egypt,'' said the Swallow. ``My friends are flying up and down the Nile,
and talking to the large lotus-flowers.''

``Swallow, Swallow, little Swallow,'' said the Prince, ``will you not stay with me for one
night, and be my messenger? The boy is so thirsty, and the mother so sad.''
\end{texte}
\source{Oscar Wilde, \og The Happy Prince \fg, \emph{The Happy Prince and Other Tales},
Londres, Nutt, 1888.}

\partiel[15 min]{Questions}
\consigne{Répondez aux deux premières questions en français, aux deux dernières en anglais.}

\begin{enumerate}[label=\arabic*.,itemsep=2.5mm,leftmargin=*]
  \item Pourquoi l'hirondelle croit-elle d'abord qu'il pleut ?
    \lignes{2}
    \corr{Parce que des gouttes tombent sur elle alors qu'il n'y a pas un nuage dans le ciel.
    Ce sont les larmes du prince, qui coulent le long de ses joues dorées.}
  \item Que demande le prince, et pourquoi ne peut-il pas le faire lui-même ?
    \lignes{2}
    \corr{D'apporter à la couturière le rubis de la poignée de son épée, pour son petit garçon
    malade. Il ne peut pas bouger : ses pieds sont fixés au socle.}
  \item \en{The Prince lived in a palace ``where sorrow is not allowed to enter''. What does
    that explain about him?}
    \lignes{3}
    \corr{He was called the Happy Prince because he had been kept from seeing anything sad,
    not because he knew what happiness was. A high wall stood round the garden and he never
    asked what lay beyond it. Now that he is set up high he sees the whole city, and cannot
    stop weeping.}
  \item \en{Find the sentence with \emph{if} in the first line. Which tense follows \emph{if},
    and what does the Swallow think a statue is for?}
    \lignes{3}
    \corr{\emph{What is the use of a statue if it cannot keep the rain off?} Après \emph{if},
    le présent. Pour l'hirondelle, une statue est un abri : elle juge tout selon son usage —
    c'est justement ce qu'elle va cesser de faire.}
\end{enumerate}

\note[Une phrase à remarquer]{\emph{I am waited for in Egypt} : passif d'un verbe à
préposition. La préposition reste accrochée au verbe et se retrouve en fin de phrase. Le
français dit \og on m'attend en Égypte \fg.}

%% ===================================================================
\section{Version}
\consigne{Traduisez le passage en français. Il suit le texte : l'hirondelle a fini par
accepter.}

\begin{version}
So the Swallow picked out the great ruby from the Prince's sword, and flew away with it in his
\gl[a beak]{beak}{un bec} over the roofs of the town.

He passed by the cathedral tower, where the white marble angels were
\gl[to sculpture]{sculptured}{sculpter}. He passed by the palace and heard the sound of
dancing. A beautiful girl came out on the balcony with her lover. ``How wonderful the stars
are,'' he said to her, ``and how wonderful is the power of love!''

``I hope my dress will be ready in time for the State-ball,'' she answered; ``I have ordered
passion-flowers to be embroidered on it; but the seamstresses are so lazy.''

He passed over the river, and saw the \gl[a lantern]{lanterns}{une lanterne} hanging to the
\gl[a mast]{masts}{un mât} of the ships. At last he came to the poor house and looked in. The
boy was \gl[to toss]{tossing}{s'agiter} \gl[feverishly]{feverishly}{fiévreusement} on his bed,
and the mother had fallen asleep, she was so tired. In he hopped, and laid the great ruby on
the table beside the woman's \gl[a thimble]{thimble}{un dé à coudre}.
\end{version}
\source{\emph{Ibid.}}

\lignes{16}

\corr{\textbf{Proposition de traduction.} L'hirondelle détacha donc le gros rubis de l'épée du
prince et s'envola avec lui dans son bec au-dessus des toits de la ville. Elle passa près de
la tour de la cathédrale, où sont sculptés les anges de marbre blanc. Elle passa près du
palais et entendit le bruit d'un bal. Une belle jeune fille sortit sur le balcon avec son
amoureux. \og Comme les étoiles sont merveilleuses, lui dit-il, et comme est merveilleuse la
puissance de l'amour ! \fg{} \og J'espère que ma robe sera prête à temps pour le bal de la
cour, répondit-elle ; j'ai commandé qu'on y brode des fleurs de la passion, mais les
couturières sont si paresseuses. \fg{} Elle passa au-dessus du fleuve et vit les lanternes
suspendues aux mâts des navires. Elle arriva enfin à la pauvre maison et regarda à
l'intérieur. Le petit garçon s'agitait fiévreusement sur son lit, et la mère s'était endormie,
tant elle était fatiguée. L'hirondelle entra d'un bond et posa le gros rubis sur la table, à
côté du dé à coudre.}

\note[Choix de traduction 1]{\emph{the Swallow\ldots his beak} : l'hirondelle est un \emph{he}
en anglais, un \og elle \fg{} en français. Choisir un genre et s'y tenir : tout le conte
repose sur ce personnage.}

\note[Choix de traduction 2]{\emph{I have ordered passion-flowers to be embroidered on it} :
infinitif passif après \emph{order}. Le français passe par une subordonnée — \og j'ai commandé
qu'on y brode \fg. La jeune fille et la couturière sont dans la même phrase sans se
rencontrer : c'est tout le conte.}

\note[Choix de traduction 3]{\emph{In he hopped} : l'adverbe en tête, puis le sujet, puis le
verbe. Tournure vive, impossible telle quelle en français, qui doit passer par un verbe et un
complément de manière — \og entra d'un bond \fg.}

%% ===================================================================
\partiel{Lexique à mémoriser}
\begin{multicols}{2}\small
\begin{itemize}[label=--,itemsep=0.8mm,leftmargin=*]
  \item \textbf{a swallow} — une hirondelle
  \item \textbf{a statue} — une statue
  \item \textbf{to weep} — pleurer ; wept, wept
  \item \textbf{a tear} — une larme \textit{(fiche 17)}
  \item \textbf{a cheek} — une joue
  \item \textbf{lead} — le plomb
  \item \textbf{a sword} — une épée \textit{(le w ne se prononce pas)}
  \item \textbf{a needle} — une aiguille
  \item \textbf{to sew} — coudre ; sewed, sewn
  \item \textbf{ill} — malade ; \textit{a fever} une fièvre
  \item \textbf{thirsty} — assoiffé ; \textit{hungry} affamé
  \item \textbf{a roof} — un toit
\end{itemize}
\end{multicols}

\end{document}
